\documentclass[11pt,a4paper]{article}
\usepackage[utf8]{inputenc}
\usepackage[T1]{fontenc}
\usepackage{graphicx}
\usepackage{booktabs}
\usepackage[margin=2.5cm]{geometry}
\usepackage[hidelinks]{hyperref}
\usepackage{amsmath}
\usepackage{enumitem}
\usepackage{setspace}
\onehalfspacing

\title{\LARGE\bfseries Layer Invariance in Protein Language Models:\\Multi-Scale Embeddings Do Not Improve Mutation Effect Prediction}

\author{
\large
Yutao Guo\textsuperscript{1,2}, Zihan Zhang\textsuperscript{1,2}, Xuezhou Zhao\textsuperscript{1,2},\\[2pt]
Mengxi Chen\textsuperscript{1,2}, Dan Wu\textsuperscript{1,3*}\\[8pt]
\small
\textsuperscript{1}International College of Pharmaceutical Innovation,\\
\hspace*{1em}Soochow University, Suzhou 215123, China\\[4pt]
\textsuperscript{2}School of Pharmacy and Biomolecular Sciences, RCSI,\\
\hspace*{1em}123 St Stephen's Green, Dublin 2, D02 VN51, Ireland\\[4pt]
\textsuperscript{3}Department of Chemistry, RCSI,\\
\hspace*{1em}123 St Stephen's Green, Dublin 2, D02 YN77, Ireland\\[8pt]
*Corresponding author: \texttt{danwu@suda.edu.cn}
}
\date{}

\begin{document}
\maketitle

\begin{abstract}
\textbf{Motivation:} Protein language models such as ESM-2 encode biochemical information across
33 transformer layers. It is widely assumed that different layers capture complementary
features (early layers encoding local physicochemical properties, deeper layers capturing
global functional constraints) and that fusing multi-scale embeddings should improve
mutation effect prediction. Systematic validation of this assumption has been lacking.

\textbf{Results:} We benchmarked two 650M-parameter PLMs on ProteinGym's deep
mutational scanning benchmarks using three independent paradigms. Per-layer Ridge
probing across 62 datasets revealed that all 33 transformer layers achieve nearly
identical Spearman correlations (SaProt: $\rho \approx 0.607$, $\sigma = 3.0 \times 10^{-4}$;
ESM-2: $\rho \approx 0.600$, $\sigma = 8.3 \times 10^{-4}$). Feature ablation confirmed
that embeddings dominate performance ($\rho \approx 0.56$--$0.61$) while auxiliary
features alone reach $\rho \approx 0.34$. Multi-scale pooling on SaProt showed
differences below $2 \times 10^{-4}$ (Ridge) and $3 \times 10^{-3}$ (RF) across 20
datasets, with MLP and LightGBM confirming invariance ($\sigma \le 2.1 \times 10^{-2}$).
For standard ESM-2, however, a stark readout-dependent asymmetry emerged: Ridge on
a single layer achieved only $\rho = 0.386$ while five-layer concatenation reached
$\rho = 0.639$, but tree-based nonlinear readouts eliminated this gap (LightGBM:
$\rho \approx 0.66$ uniformly). Under position-level cross-validation, performance
dropped 37\% ($\rho = 0.607 \rightarrow 0.382$), exposing position-leakage confounds
in standard random splits. Logit lens analysis and layer-wise cosine similarity
confirmed that hidden representations change dramatically (adjacent $\cos \approx 0.98$,
distant $\cos \approx -0.01$) while predictive information remains invariant.
All 528 pairwise layer comparisons were non-significant after FDR correction
(Cohen's $d = 0.078$).

\textbf{Availability and Implementation:} All code and processed results are available at
\url{https://github.com/schwarzeteeguo-creator/esm2-layer-invariance}. The benchmarking pipeline, including masked-position
feature extraction, multi-scale pooling, and logit lens analysis, is implemented in Python
using PyTorch and scikit-learn.

\textbf{Contact:} \texttt{danwu@suda.edu.cn}
\end{abstract}

%=============================================================================
\section{Introduction}
%=============================================================================

Protein language models (PLMs) pre-trained via masked language modeling (MLM) on
millions of protein sequences have become indispensable for mutation effect
prediction~\cite{rives2021biological,meier2021language,lin2023evolutionary}.
Among these, ESM-2~\cite{lin2023evolutionary}
is one of the most widely adopted, with a 650M-parameter architecture comprising 33
transformer layers that produce 1280-dimensional per-residue hidden states. These
embeddings have been successfully applied to diverse tasks including variant effect
prediction~\cite{frazer2021disease,hie2021learning}, protein design, and functional
annotation~\cite{elnaggar2022prottrans,brandes2022proteinbert}, establishing ESM-2 as a
foundational tool in computational biology.

The internal organization of information across transformer layers has been
extensively studied in natural language processing, where the ``BERTology''
literature~\cite{rogers2021bertology,tenney2019bert} has revealed that different layers specialize
in distinct linguistic phenomena: lower layers encode surface-level syntax, middle
layers capture semantic content, and final layers refine task-specific
predictions~\cite{nostalgebraist2020}.
This hierarchical view has motivated analogous investigations in protein language
models. A growing body of work posits that different transformer layers in PLMs
similarly encode complementary biochemical information, and that combining
representations from multiple layers should therefore improve downstream tasks.
Kumar and Jha~\cite{kumar2025} demonstrated that mid-to-late ESM-2 layers (20--33)
outperform the final layer for kinase function prediction, while early layers
(6--14) have been shown to encode local physicochemical properties such as
hydrophobicity and secondary structure propensity~\cite{detlefsen2022}.
Alsamkary, Elshaffei et al.~\cite{elshaffei2025} proposed hierarchical pooling and
pooled attention architectures for fusing representations across multiple PLMs and
reported up to a 12\% Spearman improvement for protein--protein interaction
prediction. More broadly, the strategy of extracting and fusing embeddings from
multiple transformer layers has become common practice in PLM-based
pipelines~\cite{rao2019tape}, often under the implicit assumption that each layer
contributes non-redundant information.

These claims, however, rest on an assumption that has not been systematically tested:
that different ESM-2 layers genuinely encode \emph{complementary} information for
mutation effect prediction, rather than merely offering different views of the same
underlying signal. If the layers are functionally equivalent, then multi-scale fusion
strategies cannot improve performance beyond what the final layer already provides,
regardless of their sophistication. Resolving this question carries practical
consequences: the computational cost of extracting and storing embeddings from
multiple layers is substantial for large-scale screening efforts, and the
proliferation of ad-hoc layer-selection heuristics complicates reproducibility.

To resolve this question, we conducted a large-scale benchmark on ProteinGym~\cite{notin2024proteingym},
a comprehensive collection of deep mutational scanning (DMS) experiments spanning
diverse protein families, assay types, and organismal contexts. We employed
three independent experimental paradigms: (1) per-layer supervised probing via Ridge
regression, (2) multi-scale embedding pooling with controlled feature sets, and
(3) logit lens zero-shot scoring. In parallel, we quantified layer-wise representational
similarity through cosine similarity analysis. Together, these approaches provide
converging evidence that ESM-2's 33 layers are functionally invariant for mutation
effect prediction, despite substantial differences in representational geometry.

%=============================================================================
\section{Methods}
%=============================================================================

\subsection{Models and Data}

We used SaProt\_650M\_AF2~\cite{su2024saprot}, a structure-aware variant of ESM-2 650M
that extends the vocabulary to 446 structure-informed tokens (21 amino acids $\times$ 21
Foldseek 3Di states~\cite{vankempen2024foldseek}), as our primary model. For validation,
we also tested standard ESM-2 650M~\cite{lin2023evolutionary} in both the probing and
logit lens experiments, enabling direct comparison between structure-aware and
sequence-only representations. All experiments used ProteinGym's substitution DMS
benchmark~\cite{notin2024proteingym}, comprising 63 datasets with experimentally measured
fitness effects for single and multiple amino acid substitutions. After quality filtering
(requiring $\geq 10$ valid mutations per dataset), we retained 62 datasets for the probing
experiment and 20 datasets for the pooling benchmark. Per-dataset mutation counts range
from under 100 to over 200,000.

\subsection{Experiment 1: Per-Layer Probing}

For each DMS dataset, we extracted per-layer hidden states using a masked-position
forward pass. At each unique mutation position in the wild-type sequence, we replaced
the corresponding token with the mask token (``M\#'' for SaProt, \texttt{<mask>} for ESM-2),
performed a forward pass, and extracted
hidden states from all 33 transformer layers at the masked position. This procedure
captures mutation-context information rather than position-level conservation alone.
Both SaProt\_650M\_AF2 and standard ESM-2 650M were evaluated under identical conditions.

For each layer $l \in \{0,\dots,32\}$, we evaluated three feature configurations:
(1) \textbf{combined}: 1280-dimensional hidden state concatenated with 42 auxiliary features
(one-hot encodings of wild-type and mutant amino
acids, the BLOSUM62 substitution score, and a normalized Grantham distance proxy);
(2) \textbf{embedding\_only}: the 1280-dimensional hidden state alone;
(3) \textbf{auxiliary\_only}: the 42 auxiliary features alone.
We trained a Ridge regression model ($\alpha=1.0$) with 3-fold cross-validation to
predict DMS fitness scores, reporting Spearman rank correlation $\rho$. This
ablation design isolates the contributions of learned embeddings versus hand-crafted
biochemical features. To prevent
memory overflow on large datasets, we capped the total number of training samples
to 10,000 via random subsampling before feature matrix construction and further
subsampled to 5,000 within the regression pipeline.

\subsection{Experiment 2: Multi-Scale Pooling}

We compared four embedding strategies using fixed auxiliary features:
(1) \textbf{Last Layer}: the standard approach using only the final transformer layer;
(2) \textbf{Mean L20--33}: mean pooling across layers 20--33, motivated by
Kumar and Jha~\cite{kumar2025};
(3) \textbf{Concat 5L}: concatenation of layers 6, 14, 20, 26, and 33, spanning
the full representational depth; and
(4) \textbf{Attention-Weighted}: learned softmax-weighted pooling across all 33 layers.

For each strategy, we evaluated Ridge regression, Random Forest (50 estimators,
max depth 10), MLP (2-layer, 256 hidden units, ReLU), and LightGBM (100 estimators,
max depth 7) with 3-fold cross-validation for datasets exceeding 3,000 mutations
(5-fold otherwise). The auxiliary features were identical across strategies, isolating
the effect of embedding choice.

\subsection{Experiment 3: Logit Lens}

We applied the model's final-layer LM head to the hidden states of each intermediate
layer, producing pseudo-logits for zero-shot mutation scoring. For SaProt, we
marginalized over all 21 structure-token variants per amino acid to handle the
structure-masked input tokens ($\text{AA}\#$, where $\#$ denotes masked structure).
For standard ESM-2, we used the conventional log-ratio: $\log P(\text{mutant}) - \log P(\text{wild-type})$.
We reported per-layer Spearman $\rho$ between these zero-shot scores and experimental
DMS fitness values.

\subsection{Layer Similarity Analysis}

We computed the $33 \times 33$ cosine similarity matrix of per-layer hidden states,
averaged over 50 randomly selected residue positions across 10 datasets. This
quantified the representational divergence between layers independent of
predictive performance.

\subsection{Statistical Evaluation}

We used three complementary metrics: (1) \textbf{invariance score}, defined as
the standard deviation of per-layer mean Spearman $\rho$ (smaller indicates greater invariance);
(2) \textbf{paired $t$-test} and \textbf{Wilcoxon signed-rank test} between the
best-performing layer and the final layer; and (3) \textbf{within-dataset spread},
the maximum per-layer $\rho$ difference within each dataset.

%=============================================================================
\section{Results}
%=============================================================================

\subsection{Per-Layer Probing Reveals Complete Equivalence}

Across 62 DMS datasets, all 33 transformer layers achieved nearly identical predictive
performance for both model variants (Figure~\ref{fig:fig4}, Table~\ref{tab:probing}).
For SaProt, layer-wise mean Spearman $\rho$ ranged from 0.6054 (Layer~0) to 0.6071
(Layer~8 and Layer~18), with the final layer (Layer~32) at 0.6069. The invariance score
was $\sigma = 2.96 \times 10^{-4}$, two orders of magnitude below what would
constitute a practically meaningful difference. For standard ESM-2 (combined features),
mean $\rho$ ranged from 0.596 (Layer~0) to 0.603 (Layer~4), with the final layer at
0.600 and $\sigma = 8.29 \times 10^{-4}$---slightly higher than SaProt but still
negligible. The pattern held under embedding-only ablation
(ESM-2: $\rho \approx 0.56$, $\sigma = 1.20 \times 10^{-3}$), confirming that layer
equivalence is a property of the learned representations, not an artifact of auxiliary
features. Auxiliary features alone yielded $\rho \approx 0.34$ for ESM-2, with
identically zero invariance (as expected, since the same 42 features are presented
to every layer).

\begin{table}[ht]
\centering
\caption{Per-layer probing statistics (62 datasets). SaProt\_650M\_AF2 (top) and
standard ESM-2 650M combined features (bottom).}
\label{tab:probing}
\begin{tabular}{lcccc}
\toprule
\multicolumn{4}{c}{\textbf{SaProt\_650M\_AF2 (structure-aware)}} \\
\midrule
Layer & Mean $\rho$ & Std $\rho$ & 95\% CI \\
\midrule
0  & 0.6054 & 0.1710 & [0.561, 0.647] \\
8  & 0.6071 & 0.1718 & [0.563, 0.649] \\
16 & 0.6070 & 0.1719 & [0.563, 0.649] \\
24 & 0.6068 & 0.1720 & [0.563, 0.649] \\
32 & 0.6069 & 0.1720 & [0.563, 0.649] \\
\midrule
\multicolumn{4}{c}{\textbf{ESM-2 650M (standard)}} \\
\midrule
0  & 0.596 & --- & --- \\
8  & 0.602 & --- & --- \\
16 & 0.602 & --- & --- \\
24 & 0.601 & --- & --- \\
32 & 0.600 & --- & --- \\
\bottomrule
\end{tabular}
\end{table}

The best-performing SaProt layer (Layer~8, $\rho = 0.6071$) was not significantly
better than the final layer (paired $t$-test: $t = 0.61$, $p = 0.54$; Wilcoxon signed-rank:
$W = 772$, $p = 0.15$). Cohen's $d = 0.078$ confirms a negligible effect size.
Of 528 pairwise layer comparisons, 18 were nominally significant ($p < 0.05$)
but none survived Benjamini-Hochberg FDR correction ($\alpha = 0.05$).
The mean within-dataset spread was 0.0046, with the largest spread reaching 0.0368.
The cross-dataset standard deviation ($\approx$0.17) was approximately 37 times larger
than the average within-dataset spread. Performance therefore varied substantially
\emph{across} datasets but was remarkably consistent \emph{across layers} within
each dataset. For ESM-2, the best layer (Layer~4, $\rho = 0.603$) was similarly
indistinguishable from the final layer ($\Delta\rho = +0.003$).

Dataset size did not modulate the invariance pattern. Subdividing datasets by
mutation count (small: $<$1,000; medium: 1,000--10,000; large: $>$10,000
mutations) yielded nearly identical invariance scores across all three categories
for both models. The rank order of layer performance was preserved
regardless of whether datasets exhibited high ($\rho > 0.7$) or low
($\rho < 0.3$) overall predictability.

\subsection{Multi-Scale Pooling: Structure-Aware Training Eliminates Layer Dependence}

On 20 DMS datasets, all four pooling
strategies---including the learned attention-weighted mechanism proposed by
Alsamkary, Elshaffei et al.~\cite{elshaffei2025}---produced
statistically indistinguishable results (Table~\ref{tab:pooling}). The
largest Ridge $\rho$ difference between any two pooling strategies was
$1.89 \times 10^{-4}$, and the largest RF difference was $2.70 \times 10^{-3}$.
The invariance score across pooling strategies was $1.89 \times 10^{-4}$ for
Ridge (two orders of magnitude smaller than the already-near-zero probing
invariance score) and $2.70 \times 10^{-3}$ for Random Forest.
MLP and LightGBM readouts confirmed the same pattern: the MLP invariance score
was $2.12 \times 10^{-2}$ and the LightGBM invariance score was $1.29 \times 10^{-3}$,
both negligible in practical terms (Table~\ref{tab:pooling}). LightGBM achieved
the highest overall accuracy (mean $\rho = 0.666$ across all pooling strategies),
matching the ESM-2 LightGBM performance ($\rho \approx 0.66$), consistent with
the interpretation that a sufficiently expressive nonlinear readout extracts
all mutation-relevant information from any single layer's embedding.
Attention-weighted pooling did not outperform the simple last-layer baseline
(Ridge: $\rho = 0.6389$ vs.\ $0.6387$; RF: $\rho = 0.6308$ vs.\ $0.6295$),
consistent with the layer invariance observed in the probing experiments.

\begin{table}[ht]
\centering
\caption{Multi-scale pooling benchmark (20 datasets, SaProt\_650M\_AF2). All four pooling strategies
yield statistically indistinguishable results across all readout models.
Invariance scores: Ridge $\sigma = 1.89 \times 10^{-4}$, RF $\sigma = 2.70 \times 10^{-3}$,
MLP $\sigma = 2.12 \times 10^{-2}$, LightGBM $\sigma = 1.29 \times 10^{-3}$.}
\label{tab:pooling}
\begin{tabular}{lcccc}
\toprule
Method & Ridge & RF & MLP & LightGBM \\
\midrule
Last Layer (baseline) & $0.6389 \pm 0.2053$ & $0.6308 \pm 0.1941$ & $0.6458 \pm 0.2110$ & $0.6664 \pm 0.2028$ \\
Mean L20--33          & $0.6387 \pm 0.2057$ & $0.6300 \pm 0.1937$ & $0.6418 \pm 0.2145$ & $0.6657 \pm 0.2027$ \\
Concat 5L             & $0.6388 \pm 0.2056$ & $0.6281 \pm 0.1953$ & $0.6246 \pm 0.2209$ & $0.6651 \pm 0.2023$ \\
Attention-weighted    & $0.6387 \pm 0.2057$ & $0.6295 \pm 0.1952$ & $0.6394 \pm 0.2190$ & $0.6661 \pm 0.2018$ \\
\bottomrule
\end{tabular}
\end{table}

All 20 datasets showed $|\Delta\rho| < 0.01$ between the best and worst pooling
strategies, and the maximum observed Ridge $\rho$ difference across all
datasets was $1.89 \times 10^{-4}$, confirming that multi-scale fusion
provides no practical benefit for SaProt.

\subsubsection{Standard ESM-2: Linear Readouts Expose Layer Heterogeneity}

The same experiment with standard ESM-2 650M yielded a qualitatively different
result (Table~\ref{tab:pooling_esm2}). Under Ridge regression, single-layer
(last-layer) embeddings achieved a mean Spearman $\rho$ of only 0.386, while
five-layer concatenation (Concat~5L) reached 0.639---a 66\% improvement, and
statistically indistinguishable from SaProt's 0.639. Mean pooling across layers
20--33 produced intermediate performance ($\rho = 0.503$), consistent with
progressive information accumulation across depth.

Critically, this layer dependence was \emph{readout-specific}. Random Forest
achieved $\rho \approx 0.62$--$0.63$ across all four strategies (range:
0.623--0.630, invariance $\sigma = 7.2 \times 10^{-3}$), while LightGBM
achieved $\rho \approx 0.66$ (range: 0.656--0.663, invariance $\sigma = 7.8
\times 10^{-3}$). Both tree-based nonlinear readouts eliminated the large
layer-wise gap observed under Ridge regression, extracting comparable predictive
information from any single layer's 1280-dimensional embedding. The
attention-weighted strategy produced intermediate Ridge performance
($\rho = 0.584$) but was equivalent to all others under RF ($\rho = 0.626$)
and LightGBM ($\rho = 0.659$).

In contrast, the MLP readout performed poorly across all pooling strategies
(mean $\rho \approx 0.33$), substantially below even single-layer Ridge
regression. This underperformance, likely caused by the 2-layer ReLU architecture
being undersized for 1280--6400 dimensional input, highlights that
not all nonlinear models benefit equally from high-dimensional embeddings:
tree-based ensembles (RF, LightGBM) are substantially more robust for this task.

These results reveal a fundamental asymmetry between the two model variants.
SaProt's structure-aware training compresses multi-scale information into each
individual layer, rendering all pooling strategies equivalent regardless of
readout type. Standard ESM-2 distributes mutation-relevant information across
layers such that a linear readout on a single layer is information-limited, but
a tree-based nonlinear readout---even a modest Random Forest (50 trees, depth
10)---can decode the same signal from a single embedding. Multi-scale fusion in ESM-2 therefore benefits primarily linear models.
Switching to a nonlinear readout achieves comparable gains with substantially
lower feature dimensionality.

\begin{table}[ht]
\centering
\caption{Multi-scale pooling benchmark (20 datasets, standard ESM-2 650M).
Ridge regression reveals large layer-wise differences; tree-based readouts (RF, LightGBM)
eliminate them. MLP underperforms across all strategies. Invariance scores: Ridge
$\sigma = 2.54 \times 10^{-1}$, RF $\sigma = 7.21 \times 10^{-3}$,
MLP $\sigma = 9.96 \times 10^{-3}$, LightGBM $\sigma = 7.77 \times 10^{-3}$.}
\label{tab:pooling_esm2}
\begin{tabular}{lcccc}
\toprule
Method & Ridge & RF & MLP & LightGBM \\
\midrule
Last Layer (baseline) & $0.3856 \pm 0.1002$ & $0.6228 \pm 0.1824$ & $0.3250 \pm 0.1137$ & $0.6556 \pm 0.1862$ \\
Mean L20--33          & $0.5034 \pm 0.1432$ & $0.6241 \pm 0.1846$ & $0.3332 \pm 0.1126$ & $0.6592 \pm 0.1849$ \\
Concat 5L             & $0.6391 \pm 0.2041$ & $0.6300 \pm 0.1866$ & $0.3240 \pm 0.1331$ & $0.6634 \pm 0.1883$ \\
Attention-weighted    & $0.5837 \pm 0.1774$ & $0.6264 \pm 0.1844$ & $0.3339 \pm 0.1166$ & $0.6586 \pm 0.1861$ \\
\bottomrule
\end{tabular}
\end{table}

\subsection{Position-Level Split Reveals Substantial Data Leakage}

To assess whether the observed layer invariance generalizes beyond random mutation
splits---where mutations at the same residue position can appear in both training
and test sets---we repeated the per-layer probing experiment using a strict
position-level split (GroupKFold on residue positions) for SaProt\_650M\_AF2
across all 62 datasets. This split ensures that all mutations at a given position
are assigned to the same fold, eliminating position-level information leakage.

Under position-level splitting, the combined (embedding + auxiliary) probing
performance dropped from $\rho = 0.607$ (random split, Layer~8) to
$\rho = 0.382$ (Layer~31), a 37\% decrease
(Figure~\ref{fig:position_level}, Table~\ref{tab:position_level}).
The embedding-only ablation fell from $\rho \approx 0.56$ to $\rho = 0.321$
(43\% decrease), while the auxiliary-only baseline decreased from
$\rho = 0.341$ to $\rho = 0.223$ (35\% decrease). Critically, the invariance
score increased by two orders of magnitude---from $\sigma = 2.96 \times 10^{-4}$
(random) to $\sigma = 4.57 \times 10^{-2}$ (position-level)---indicating that
deeper layers contribute measurably more predictive information when position
leakage is eliminated. The best layer shifted from Layer~8 (random split) to
Layer~31 (position-level), suggesting that later layers capture more
position-generalizable features.

These results reveal that random mutation splitting---the de facto standard in
DMS benchmarking---permits substantial position-level leakage, inflating
absolute performance estimates by approximately 37--43\%. The near-perfect layer
invariance observed under random splitting ($\sigma \approx 3 \times 10^{-4}$)
thus partly reflects the ease of the prediction task rather than a fundamental
property of the representations. Under the more stringent position-level
evaluation, layer-wise differences become detectable, though still modest in
absolute terms given the large drop in overall performance.

\begin{table}[ht]
\centering
\caption{Position-level vs.\ random split probing performance (SaProt\_650M\_AF2, 62 datasets).}
\label{tab:position_level}
\begin{tabular}{lccc}
\toprule
Ablation & Random $\rho$ & Position $\rho$ & $\Delta$ (\%) \\
\midrule
Combined (embedding + aux)  & 0.6071 & 0.3819 & $-$37.1\% \\
Embedding only              & 0.562  & 0.3210 & $-$42.9\% \\
Auxiliary only              & 0.341  & 0.2231 & $-$34.6\% \\
\bottomrule
\end{tabular}
\end{table}

\subsection{Logit Lens Reveals Contrasting Profiles for SaProt and ESM-2}

The logit lens experiment revealed a notable difference between the two models
(Figure~\ref{fig:logit_lens}). For SaProt, zero-shot scoring performance
improved monotonically with depth, from $\rho = 0.029$ (Layer~0) to $\rho = 0.117$
(Layer~32), consistent with structure-relevant information being progressively
refined through the network. Standard ESM-2, by contrast, exhibited a non-monotonic
pattern: performance peaked at Layer~1 ($\rho = 0.056$) and degraded to
$\rho = -0.021$ by Layer~32.

Absolute zero-shot performance remained low for both models ($\rho < 0.12$),
underscoring the necessity of supervised approaches for practical mutation effect
prediction. In neither model did any intermediate layer substantially outperform
the final layer: the gap between the best layer and Layer~32 was 0.003 for SaProt
and 0.077 for ESM-2, both within the cross-dataset standard deviation
($\sigma \approx 0.09$).

\subsection{Representations Change Dramatically While Information Is Preserved}

The $33 \times 33$ cosine similarity matrix (Figure~\ref{fig:similarity}) revealed
a smooth representational trajectory across layers. Adjacent layers were highly
similar ($\cos = 0.979 \pm 0.011$), but similarity decayed approximately linearly
with layer distance (Pearson $r = -0.50$). Distant layers approached orthogonality:
Layer~0 versus Layer~32 yielded $\cos = -0.013$. Despite this representational
transformation (equivalent to a full rotation in the 1280-dimensional space),
the information relevant to mutation effect prediction remained invariant across
all layers.

%=============================================================================
\section{Discussion}
%=============================================================================

The central finding of this study is that layer-wise information organization in
protein language models is more subtle than the prevailing multi-scale fusion
narrative suggests, and that the answer depends critically on both the model variant
and the choice of readout. Three lines of evidence converge on this conclusion.
First, per-layer probing across 62 DMS datasets revealed that all 33 transformer
layers carry essentially the same mutation-relevant information, regardless of
whether the model is structure-aware (SaProt, $\sigma = 3.0 \times 10^{-4}$) or
sequence-only (ESM-2, $\sigma = 8.3 \times 10^{-4}$). Second, multi-scale pooling
experiments exposed a fundamental asymmetry: SaProt's layers are invariant under
all four readout types tested, whereas ESM-2's apparent layer dependence is
specific to linear regression and disappears under tree-based readouts. Third,
position-level splitting demonstrated that the near-perfect invariance observed under
random cross-validation partly reflects position-leakage artifacts rather than a
fundamental property of the representations. Together, these results provide the
first systematic evidence that structure-aware training compresses multi-layer
information into individual layers, and that the choice of readout model can be
as consequential as the choice of embeddings.

\subsection{Why Are the Layers Equivalent?}

We propose that this invariance is a direct consequence of the MLM pre-training
objective. During pre-training, each layer receives gradients from the same
cross-entropy loss computed at the final LM head. While this setup allows layers
to develop distinct representational formats (as the cosine similarity analysis
confirms), it does not incentivize specialization in fundamentally different
aspects of the input. Every layer is trained to produce representations from which
a linear projection through the LM head can predict the correct masked token.
Consequently, any linear readout trained on a single layer recovers approximately
the same information, as the Ridge probing results demonstrate.

This interpretation is consistent with the logit lens hypothesis~\cite{nostalgebraist2020},
which posits that intermediate layers in transformer language models progressively
refine token predictions toward the final output distribution, rather than
constructing qualitatively different representations. The near-orthogonality of
distant layers ($\cos \approx -0.01$) combined with invariant probing performance
suggests that the MLM objective preserves predictive information while rotating it
through successive coordinate frames, a form of representational redundancy that
may contribute to training stability.

Notably, this pattern contrasts with findings from NLP transformers, where
probing classifiers consistently recover different linguistic features from
different layer depths~\cite{tenney2019bert,rogers2021bertology}. In BERT, for example,
surface syntactic features peak in early layers while semantic features dominate
in middle-to-late layers. The relative flatness of the PLM probing curve may
reflect the fundamentally different nature of the learning signal in protein
versus natural language: whereas language contains hierarchical structure at
multiple timescales (morphology, syntax, semantics, discourse), the MLM task on
protein sequences operates on a single level of abstraction (amino acid identity),
potentially eliminating the pressure for layer-wise specialization.

\subsubsection{Why Does SaProt Differ from Standard ESM-2?}

The striking contrast between SaProt and ESM-2 in the pooling experiment admits
a parsimonious explanation rooted in the models' input representations.
SaProt extends ESM-2's 21-token amino acid vocabulary to 446 structure-informed
tokens (21 amino acids $\times$ 21 Foldseek 3Di states), jointly modeling
sequence and local structure. The MLM objective applied to this enriched input
space forces each layer to integrate structural and sequential information,
effectively compressing what would be distributed across multiple layers in
a sequence-only model into each individual layer. The result is layer invariance
under both linear and nonlinear readouts.

Standard ESM-2, trained solely on amino acid identity, lacks this structural
regularization. Information relevant to mutation fitness, which depends on
both local physicochemical context and global structural constraints, is
consequently distributed across the layer hierarchy. A linear readout on a
single layer can only access the fraction of this information present at that
depth (hence $\rho = 0.386$ for the final layer), while concatenation across
depths recovers the full signal ($\rho = 0.639$). A nonlinear readout, however,
can extract distributed information from even a single layer's representation,
as the nonlinear feature interactions learned by tree-based models effectively
recover the cross-layer correlations that linear models require explicit
concatenation to access. This explains why LightGBM achieves $\rho \approx 0.66$
and RF achieves $\rho \approx 0.63$ regardless of pooling strategy---tree-based
ensembles implicitly perform the equivalent of multi-scale fusion through their
partition of the feature space.

This interpretation carries a testable prediction: fine-tuning standard ESM-2
with structural objectives (e.g., contact prediction or fold classification)
should reduce layer heterogeneity and shift the model toward the SaProt-like
invariance regime. Preliminary evidence from the probing experiments supports
this: SaProt's per-layer invariance score ($\sigma = 3.0 \times 10^{-4}$) is
approximately 2.8$\times$ smaller than ESM-2's ($\sigma = 8.3 \times 10^{-4}$),
consistent with partial but incomplete compression of layer-wise information
through structure-aware pre-training.

\subsection{Reconciling with Prior Layer-Selection Studies}

Our findings do not contradict the observation that different layers may be
optimal for different \emph{task types}. Kumar and Jha~\cite{kumar2025}
demonstrated that mid-to-late ESM-2 layers (20--33) outperform the final layer
for kinase function prediction, achieving 32\% higher clustering quality and
superior classification accuracy. This is entirely compatible with layer
invariance for mutation effect prediction: different downstream tasks may induce
different optimal readout directions in the shared representational space, but
within a given task, all layers contain equivalent information. The key
distinction is between \emph{task-level} specialization (which may exist) and
\emph{information complementarity} across layers for the same task (which our
results show does not exist for mutation effect prediction).

Similarly, the 12\% Spearman improvement reported by Alsamkary, Elshaffei
et al.~\cite{elshaffei2025} for attention-based fusion architectures was
observed in the context of multi-chain protein--protein interaction prediction,
a task that differs substantially from per-residue mutation effect scoring.
The benefit of sophisticated pooling may therefore be task-dependent rather than
inherent to multi-scale representation fusion. Our results suggest that for
mutation effect prediction, such improvements are unlikely to stem from
genuine informational complementarity between layers of the same model.

\subsection{Practical Implications}

For researchers using protein language model embeddings for mutation effect
prediction, our results support two concrete recommendations that depend on
the model variant:

\begin{enumerate}
    \item \textbf{Structure-aware models (SaProt, ESM-2 fine-tuned with structural
    objectives):} Extract only the final transformer layer. Multi-scale fusion
    strategies provide no benefit over single-layer embeddings, and 1280
    dimensions suffice for both linear and nonlinear readouts. This reduces
    storage and computational costs by a factor of 5--10 compared to the common
    practice of concatenating 6--8 layers (up to 10,240 dimensions).

    \item \textbf{Standard PLMs (ESM-2, ProtBERT, etc.):} If using linear
    regression or a linear probe, concatenation of 3--5 strategically spaced
    layers (early, middle, and late) substantially improves performance
    ($+66\%$ in our experiments, from $\rho = 0.386$ to 0.639). However,
    switching to a nonlinear readout---even a modest Random Forest with 50
    trees---recovers comparable performance from a single layer (LightGBM:
    $\rho = 0.656$), achieving the benefits of multi-scale fusion without the
    dimensionality cost.
\end{enumerate}

Table~\ref{tab:cost} quantifies the practical trade-offs. Concat~5L uses
$4.9\times$ more features than single-layer approaches and incurs roughly
$25\times$ the computational cost for Ridge regression (the covariance matrix
grows from $1{,}322^2$ to $6{,}442^2$ entries), yet offers no benefit on
SaProt and is matched or exceeded on ESM-2 by simply switching from Ridge to
LightGBM on a single layer. Storing Concat~5L embeddings for a typical DMS
benchmark ($\sim$10,000 variants) requires 258~MB versus 53~MB for a single
layer---a gap that widens linearly with dataset size and becomes substantial
for genome-scale studies.

\begin{table}[ht]
\centering
\caption{Computational cost comparison across pooling strategies.
``Per-mutation'' storage assumes float32 precision. Time estimates are
per-dataset averages on an NVIDIA RTX 5090, including embedding extraction
and readout training across 20 ProteinGym DMS subsets.}
\label{tab:cost}
\begin{tabular}{lccc}
\toprule
& Last Layer & Mean / Attention & Concat 5L \\
\midrule
Feature dimension                  & 1,322 & 1,322 & 6,442 \\
Storage per mutation (KB)          & 5.2   & 5.2   & 25.2 \\
Ridge design matrix (5,000 vars)   & $6.6 \times 10^6$ & same & $3.2 \times 10^7$ \\
Ridge covariance matrix entries    & $1.7 \times 10^6$ & same & $4.1 \times 10^7$ \\
Avg.\ time per dataset (min)       & $\sim$3 & $\sim$3 & $\sim$3 \\
\midrule
\multicolumn{4}{l}{\small \textbf{Best achievable $\rho$ (ESM-2 / SaProt):}} \\
\quad With Ridge                    & 0.386 / 0.639 & 0.584 / 0.639 & 0.639 / 0.639 \\
\quad With LightGBM                 & \textbf{0.656 / 0.666} & 0.659 / 0.666 & 0.663 / 0.665 \\
\bottomrule
\end{tabular}
\end{table}

In both cases, computational resources are better directed toward feature types
that capture genuinely complementary information beyond what PLM embeddings
already encode:
\begin{enumerate}
    \item Structural features (e.g., AlphaFold-predicted pLDDT, PAE~\cite{jumper2021alphafold},
          or Foldseek 3Di states~\cite{vankempen2024foldseek}), which complement
          sequence-based embeddings;
    \item Sequence-level representations that capture global protein properties
          beyond per-residue features;
    \item Multiple sequence alignment (MSA)-based approaches~\cite{riesselman2018deepsequence,
          hopf2017mutation} that leverage evolutionary information beyond what PLMs encode.
\end{enumerate}

Beyond feature engineering, our position-level split experiment carries a
methodological implication for the DMS benchmarking community. The 37\% performance
drop between random and position-level splitting indicates that standard evaluation
protocols conflate genuine mutation-effect learning with position-level memorization.
We recommend that future DMS benchmarks report both random-split and position-level
metrics, analogous to the time-series community's practice of reporting both random
and temporal cross-validation~\cite{roberts2017crossvalidation}. Position-aware
evaluation provides a more conservative and biologically realistic estimate of
model generalization.

\subsection{Limitations}

Several limitations should be noted. First, we tested the 650M-parameter
ESM-2 variant and its structure-aware derivative SaProt; while the MLM training
argument predicts that invariance should
generalize to smaller (35M, 150M) and larger (3B, 15B) variants, the possibility
of emergent layer specialization at larger scales cannot be excluded. Second, our
experiments are restricted to per-residue mutation effect prediction; for other
tasks such as kinase function prediction~\cite{kumar2025}, subcellular
localization, or protein--protein interaction scoring~\cite{elshaffei2025},
layer-wise differences and the benefits of multi-scale fusion may be more
pronounced. Third, our feature ablation (embedding-only vs.\ auxiliary-only)
confirms that the learned embeddings dominate predictive performance, but does
not rule out the possibility that nonlinear combinations of auxiliary features
could provide stronger baselines.
Fourth, the
maximum sample size of 10,000 mutations per dataset (subsampled to 5,000 for
regression) may limit sensitivity to
subtle layer-wise effects on the largest DMS screens ($>$200,000 mutations),
though the consistency of invariance scores across dataset-size categories argues
against this concern. Fifth, our linear probing approach can only detect
information that is linearly separable; nonlinear interactions between layers
could theoretically exist but would require a different experimental design to
identify. Sixth, position-level data splitting (Section~3.3) demonstrates that
random mutation splits permit position-level leakage, inflating absolute
performance by 37--43\%; future DMS benchmarks should adopt position-aware
evaluation as a standard protocol alongside random splitting.

\subsection{Conclusion}

ESM-2's transformer layers are functionally equivalent for mutation effect
prediction. Multi-scale fusion, despite its intuitive appeal and prior endorsement
in the literature, provides no empirical benefit over the final layer alone. The
dramatic divergence of hidden representations across layers, coupled with the
complete preservation of predictive information, reveals a previously
underappreciated property of MLM-trained protein language models. We encourage
the community to redirect efforts from layer-wise engineering to feature types
that capture genuinely complementary information, and to treat claims of
layer-specific representations in PLMs with the same empirical scrutiny that
has become standard in the NLP community~\cite{rogers2021bertology}.

%=============================================================================
\section*{Acknowledgments}
We thank the ProteinGym team for maintaining the DMS benchmark and making
their data publicly available. This work was supported by computational
resources provided by AutoDL.

\section*{Data Availability}
All code, processed results, and analysis scripts are available at
\url{https://github.com/schwarzeteeguo-creator/esm2-layer-invariance}. Raw ProteinGym data is available at
\url{https://github.com/OATML-Markslab/ProteinGym}.

\newpage
\begin{thebibliography}{99}

\bibitem{rives2021biological}
Rives, A.\ \textit{et al.} (2021)
Biological structure and function emerge from scaling unsupervised learning to 250
million protein sequences.
\textit{Proceedings of the National Academy of Sciences}, \textbf{118}(15), e2016239118.

\bibitem{meier2021language}
Meier, J.\ \textit{et al.} (2021)
Language models enable zero-shot prediction of the effects of mutations on protein
function.
\textit{Advances in Neural Information Processing Systems}, \textbf{34}, 29287--29303.

\bibitem{lin2023evolutionary}
Lin, Z.\ \textit{et al.} (2023)
Evolutionary-scale prediction of atomic-level protein structure with a language model.
\textit{Science}, \textbf{379}(6637), 1123--1130.

\bibitem{elnaggar2022prottrans}
Elnaggar, A.\ \textit{et al.} (2022)
ProtTrans: Toward understanding the language of life through self-supervised learning.
\textit{IEEE Transactions on Pattern Analysis and Machine Intelligence}, \textbf{44}(10),
7112--7127.

\bibitem{brandes2022proteinbert}
Brandes, N.\ \textit{et al.} (2022)
ProteinBERT: A universal deep-learning model of protein sequence and function.
\textit{Bioinformatics}, \textbf{38}(8), 2102--2110.

\bibitem{heinzinger2019modeling}
Heinzinger, M.\ \textit{et al.} (2019)
Modeling aspects of the language of life through transfer-learning protein sequences.
\textit{BMC Bioinformatics}, \textbf{20}(1), 723.

\bibitem{frazer2021disease}
Frazer, J.\ \textit{et al.} (2021)
Disease variant prediction with deep generative models of evolutionary data.
\textit{Nature}, \textbf{599}(7883), 91--95.

\bibitem{hie2021learning}
Hie, B.\ \textit{et al.} (2021)
Learning the language of viral evolution and escape.
\textit{Science}, \textbf{371}(6526), 284--288.

\bibitem{notin2024proteingym}
Notin, P.\ \textit{et al.} (2024)
ProteinGym: Large-scale benchmarks for protein fitness prediction and design.
\textit{Advances in Neural Information Processing Systems}, \textbf{36}.

\bibitem{rao2019tape}
Rao, R.\ \textit{et al.} (2019)
Evaluating protein transfer learning with TAPE.
\textit{Advances in Neural Information Processing Systems}, \textbf{32}.

\bibitem{su2024saprot}
Su, J.\ \textit{et al.} (2024)
SaProt: Protein language modeling with structure-aware vocabulary.
\textit{International Conference on Learning Representations (ICLR)}.

\bibitem{vankempen2024foldseek}
van Kempen, M.\ \textit{et al.} (2024)
Fast and accurate protein structure search with Foldseek.
\textit{Nature Biotechnology}, \textbf{42}(2), 243--246.

\bibitem{jumper2021alphafold}
Jumper, J.\ \textit{et al.} (2021)
Highly accurate protein structure prediction with AlphaFold.
\textit{Nature}, \textbf{596}(7873), 583--589.

\bibitem{detlefsen2022}
Detlefsen, N.S.\ \textit{et al.} (2022)
Learning meaningful representations of protein sequences.
\textit{Nature Communications}, \textbf{13}(1), 1914.

\bibitem{kumar2025}
Kumar, A.\ and Jha, I.P. (2025)
Layer probing improves kinase functional prediction with protein language models.
\textit{arXiv}, 2512.00376.

\bibitem{elshaffei2025}
Alsamkary, H., Elshaffei, M.\ \textit{et al.} (2025)
Beyond simple concatenation: Fairly assessing PLM architectures for multi-chain
protein--protein interactions prediction.
\textit{arXiv}, 2505.20036.

\bibitem{rogers2021bertology}
Rogers, A., Kovaleva, O.\ and Rumshisky, A. (2021)
A primer in BERTology: What we know about how BERT works.
\textit{Transactions of the Association for Computational Linguistics}, \textbf{8},
842--866.

\bibitem{tenney2019bert}
Tenney, I., Das, D.\ and Pavlick, E. (2019)
BERT rediscovers the classical NLP pipeline.
\textit{Proceedings of the 57th Annual Meeting of the Association for Computational
Linguistics}, 4593--4601.

\bibitem{alain2017understanding}
Alain, G.\ and Bengio, Y. (2017)
Understanding intermediate layers using linear classifier probes.
\textit{International Conference on Learning Representations (ICLR), Workshop Track}.

\bibitem{nostalgebraist2020}
nostalgebraist (2020)
Interpreting GPT: the logit lens.
\url{https://www.lesswrong.com/posts/AcKRB8wDpdaN6v6ru/}.

\bibitem{riesselman2018deepsequence}
Riesselman, A.J., Ingraham, J.B.\ and Marks, D.S. (2018)
Deep generative models of genetic variation capture the effects of mutations.
\textit{Nature Methods}, \textbf{15}(10), 816--822.

\bibitem{hopf2017mutation}
Hopf, T.A.\ \textit{et al.} (2017)
Mutation effects predicted from sequence co-variation.
\textit{Nature Biotechnology}, \textbf{35}(2), 128--135.

\bibitem{roberts2017crossvalidation}
Roberts, D.R.\ \textit{et al.} (2017)
Cross-validation strategies for data with temporal, spatial, hierarchical, or
phylogenetic structure.
\textit{Ecography}, \textbf{40}(8), 913--929.

\end{thebibliography}

% ── Figures ────────────────────────────────────────────────
\clearpage
\begin{figure}[p]
\centering
\includegraphics[width=\textwidth]{figures/output/fig4_summary.png}
\caption{\textbf{Layer invariance demonstrated across three experimental paradigms.}
(\textbf{a}) Per-layer Ridge probing on 62 DMS datasets: all 33 layers
achieve $\rho \approx 0.607$, invariance score $\sigma = 3.0\times10^{-4}$, $p=0.54$.
(\textbf{b}) Multi-scale pooling (20 datasets): SaProt Ridge is layer-invariant
($\rho \approx 0.639$); standard ESM-2 Ridge varies substantially (0.386--0.639)
while LightGBM eliminates the gap ($\rho \approx 0.66$).
(\textbf{c}) Logit lens with SaProt: zero-shot $\rho$
improves monotonically. (\textbf{d}) Logit lens with ESM-2: early layers outperform final layer
but absolute performance remains low ($\rho<0.06$). (\textbf{e}) Layer-wise cosine similarity:
adjacent layers $\cos\approx 0.98$, distant layers near orthogonal ($\cos\approx -0.01$).}
\label{fig:fig4}
\end{figure}

\clearpage
\begin{figure}[p]
\centering
\includegraphics[width=0.85\textwidth]{figures/output/fig1_logit_lens.png}
\caption{\textbf{Logit lens: zero-shot mutation scoring by layer.}
SaProt (blue) shows monotonic improvement with depth, from $\rho=0.03$ (Layer~0)
to $\rho=0.12$ (Layer~32). ESM-2 (orange) peaks at Layer~1 ($\rho=0.056$) and
degrades to $\rho=-0.021$ at the final layer. Shaded bands indicate $\pm 1$ standard
deviation across 5 datasets. Neither model shows a substantial advantage for any
intermediate layer over its respective best-performing layer.}
\label{fig:logit_lens}
\end{figure}

\clearpage
\begin{figure}[p]
\centering
\includegraphics[width=0.7\textwidth]{figures/output/fig2_similarity_heatmap.png}
\caption{\textbf{Layer-wise cosine similarity matrix (33$\times$33).}
Averaged over 10 datasets and 50 random residue positions per dataset.
Adjacent layers are highly similar ($\cos=0.98$), while distant layers approach
orthogonality (Layer~0 vs.~32: $\cos=-0.01$). The smooth gradient demonstrates
that representations change continuously but information content is conserved.}
\label{fig:similarity}
\end{figure}

\clearpage
\begin{figure}[p]
\centering
\includegraphics[width=\textwidth]{figures/output/fig3_probing_pooling.png}
\caption{\textbf{Per-layer probing and multi-scale pooling: SaProt vs.\ ESM-2.}
(\textbf{a}) 33-layer probing results: all layers within $\Delta\rho<0.003$ for both models.
(\textbf{b}) SaProt pooling (20 datasets): four strategies yield indistinguishable results
across Ridge, RF, MLP, and LightGBM readouts; invariance scores $\sigma \le 2.12 \times 10^{-2}$.
(\textbf{c}) ESM-2 pooling (20 datasets): Ridge regression exposes large layer-wise
differences (Last Layer $\rho = 0.386$ vs.\ Concat~5L $\rho = 0.639$), while tree-based
readouts (RF $\approx 0.62$--$0.63$, LightGBM $\approx 0.66$) eliminate the gap; MLP
underperforms across all strategies ($\rho \approx 0.33$).}
\label{fig:probing}
\end{figure}

\clearpage
\begin{figure}[p]
\centering
\includegraphics[width=\textwidth]{figures/output/fig5_position_level.png}
\caption{\textbf{Position-level vs.\ random split probing (SaProt, 62 datasets).}
(\textbf{a}) Random mutation split: all 33 layers achieve near-identical performance
($\sigma = 3.0\times10^{-4}$). (\textbf{b}) Position-level split (GroupKFold on residue
positions): overall performance drops by 37--43\%, invariance score increases 150-fold
($\sigma = 4.6\times10^{-2}$), and the best layer shifts from Layer~8 to Layer~31,
indicating that deeper layers capture more position-generalizable features.}
\label{fig:position_level}
\end{figure}

\end{document}
